\documentclass[11pt,a4paper]{article}
\usepackage[margin=25mm,headheight=14pt]{geometry}
\usepackage[T1]{fontenc}
\usepackage[utf8]{inputenc}
\usepackage{lmodern,amsmath,amssymb}
\usepackage{xcolor,graphicx,booktabs,longtable,array}
\usepackage{microtype,enumitem,fancyhdr,titlesec}
\usepackage{hyperref}
\definecolor{accent}{HTML}{163B55}
\hypersetup{colorlinks=true,linkcolor=accent,urlcolor=accent,
 pdftitle={Local matrix-mass bounds do not imply Stahl--Totik regularity},pdfauthor={Matthew J. Colbrook},
 pdfsubject={AIM 1035: mathematical result and proof}}
\urlstyle{same}
\Urlmuskip=0mu plus 1mu
\setlength{\parindent}{0pt}
\setlength{\parskip}{0.35em}
\setlength{\emergencystretch}{3em}
\setlist{itemsep=0.25em,topsep=0.3em}
\providecommand{\tightlist}{\setlength{\itemsep}{0pt}\setlength{\parskip}{0pt}}
\titleformat{\section}{\normalfont\large\bfseries\color{accent}}{\thesection}{0.7em}{}
\titleformat{\subsection}{\normalfont\normalsize\bfseries}{\thesubsection}{0.6em}{}
\pagestyle{fancy}
\fancyhf{}
\fancyhead[L]{\small\sffamily AIM 1035}
\fancyhead[R]{\small\sffamily Matthew J. Colbrook}
\fancyfoot[C]{\thepage}
\renewcommand{\headrulewidth}{0.3pt}
\setcounter{secnumdepth}{0}
\allowdisplaybreaks[2]
\title{\textbf{Local matrix-mass bounds do not imply Stahl--Totik regularity}}
\author{Matthew J. Colbrook\\[3pt]
{\small Department of Applied Mathematics and Theoretical Physics}\\
{\small University of Cambridge, United Kingdom}\\
{\small\href{mailto:mjc249@maths.cam.ac.uk}{\texttt{mjc249@maths.cam.ac.uk}}}}
\date{3 November 2026}
\begin{document}
\maketitle
\begin{abstract}
We construct a normalized nontrivial positive-semidefinite matrix measure with a uniform polynomial local-mass lower bound whose block-Jacobi determinant regularity limit nevertheless vanishes. The rank-one density has no polynomial null vector. This gives a counterexample to the precise matrix local-mass criterion in AIM 1035.
\end{abstract}
\noindent\textbf{Repository target.}
This manuscript concerns AIM problem 1035 at revision
\texttt{3fa50884d129}. Its scope is stated below.
\medskip

\hypertarget{1035-status-and-scope}{%
\section{Status and scope}\label{1035-status-and-scope}}

This counterexample disproves AIM 1035 and Conjecture~5.9 in the work of D. Damanik,
A. Pushnitski and B. Simon, \emph{The Analytic Theory of Matrix
Orthogonal Polynomials}, Surveys in Approximation Theory \textbf{4}
(2008), 1--85 \cite{DPS08}. It uses the positive-semidefinite matrix measures allowed
by that paper. The polynomial moment form is strictly positive, as
required by \cite[Section~1.2 and Lemma~2.1]{DPS08}. The density has rank one, so a variant requiring
a positive-definite density almost everywhere is not addressed.

Primary published text:
\url{https://www.math.auckland.ac.nz/sat/papers/11/11.pdf}. The
relevant source equations are (2.34) and (5.20); the needed facts are
proved below.

\hypertarget{1035-theorem}{%
\section{Theorem}\label{1035-theorem}}

There exists a normalized, nontrivial, positive \(2\times2\) matrix
measure \(\mu\), supported on \(E=[-1,1]\), whose associated block
Jacobi matrix has essential spectrum exactly \(E\), such that

\[
\mu([x-m^{-1},x+m^{-1}])\succeq c m^{-3} I_2
\qquad(x\in E,\ m\ge1)
\]

for a constant \(c>0\), but whose invertible off-diagonal Jacobi blocks
satisfy

\[
\lim_{n\to\infty}|\det(A_1\cdots A_n)|^{1/n}=0.
\]

Since \(\operatorname{Cap}([-1,1])^2=1/4\), this measure is not regular.
Its mass condition is stronger than the conjectured condition: for every
\(\eta>0\), one has \(E\subset S_{m,\eta}\) for all sufficiently large
\(m\).

\hypertarget{1035-construction-and-nondegeneracy}{%
\section{1. Construction and
nondegeneracy}\label{1035-construction-and-nondegeneracy}}

Set

\[
v(x)=\begin{pmatrix}1\\e^x\end{pmatrix},\qquad
d\nu(x)=v(x)v(x)^T\,\mathbf1_{[-1,1]}(x)\,dx.
\]

Let \(M=\nu(E)\), \(C=M^{-1/2}\), and \(d\mu=C\,d\nu\,C\). In
particular, \(\mu(E)=I_2\). Explicitly,

\[
M=\begin{pmatrix}
2&e-e^{-1}\\
e-e^{-1}&(e^2-e^{-2})/2
\end{pmatrix},\qquad
\det M=2(1-e^{-2})>0.
\]

For a vector polynomial \(f=(a,b)^T\),

\[
\|f\|_\nu^2=\int_{-1}^1|a(x)+e^x b(x)|^2\,dx.
\]

This norm is zero only if \(a=b=0\). Indeed, vanishing of the integral
forces the entire function \(a(z)+e^z b(z)\) to vanish identically. If
\(b\ne0\), its exponential term dominates the polynomial \(a\) on the
positive real axis, a contradiction. The invertible constant change
\(f\mapsto Cf\) gives the same nondegeneracy for \(\mu\).

Consequently every finite vector-polynomial Gram matrix is positive
definite. Block Gram--Schmidt therefore produces the monic matrix
orthogonal polynomials of all degrees and an associated block Jacobi
matrix with all off-diagonal blocks invertible. The rank-one density is
compatible with strict positivity of the polynomial moment form: there
is no nonzero polynomial vector lying in its pointwise nullspace.

\hypertarget{1035-the-associated-essential-spectrum}{%
\section{2. The associated essential
spectrum}\label{1035-the-associated-essential-spectrum}}

Define, initially on vector polynomials,

\[
(Uf)(x)=v(x)^T C f(x).
\]

By construction, \(U\) is an isometry from the polynomial-vector inner
product for \(\mu\) into \(L^2([-1,1],dx)\). Its range contains every
scalar polynomial: for a scalar polynomial \(r\), take
\(f=C^{-1}(r,0)^T\). Scalar polynomials are dense in this scalar \(L^2\)
space. Thus \(U\) extends to a unitary map from the completion of the
polynomial vectors onto \(L^2([-1,1],dx)\).

Moreover, \(U(xf)=xUf\). Hence the block Jacobi operator obtained from
the orthonormal polynomial-vector basis is unitarily equivalent to
multiplication by \(x\) on scalar \(L^2([-1,1],dx)\). This is a bounded
self-adjoint operator of norm one, with spectrum and essential spectrum
both \([-1,1]\). (Its spectrum is purely absolutely continuous of
multiplicity one; the conjecture does not assume maximal matrix spectral
multiplicity.)

\hypertarget{1035-uniform-local-matrix-mass-lower-bound}{%
\section{3. Uniform local matrix-mass lower
bound}\label{1035-uniform-local-matrix-mass-lower-bound}}

Let \(I\subset[-1,1]\) be an interval of length \(L>0\). Write

\[
G_I=\nu(I)=\begin{pmatrix}
L&\int_I e^t\,dt\\
\int_I e^t\,dt&\int_I e^{2t}\,dt
\end{pmatrix}.
\]

The determinant identity

\[
\det G_I
=\frac12\int_I\int_I(e^s-e^t)^2\,ds\,dt
\]

and the mean-value bound \(|e^s-e^t|\ge e^{-1}|s-t|\) imply

\[
\det G_I\ge\frac{e^{-2}}2\int_I\int_I(s-t)^2\,ds\,dt
=\frac{e^{-2}}{12}L^4.
\]

Also \(\operatorname{tr}G_I\le(1+e^2)L\). Since \(G_I\) is positive
definite,

\[
\lambda_{\min}(G_I)
=\frac{\det G_I}{\lambda_{\max}(G_I)}
\ge\frac{e^{-2}}{12(1+e^2)}L^3.
\]

For \(x\in[-1,1]\) and \(m\ge1\), the truncated interval
\(I=[x-m^{-1},x+m^{-1}]\cap[-1,1]\) has length at least \(m^{-1}\),
including when \(x\) is an endpoint. Consequently,

\[
\mu([x-m^{-1},x+m^{-1}])=CG_I C
\succeq \frac{e^{-2}}{12(1+e^2)\lambda_{\max}(M)}m^{-3}I_2.
\]

Take the positive constant on the right as \(c\). Because
\(c m^{-3}\ge e^{-\eta m}\) for every sufficiently large \(m\), for each
fixed \(\eta>0\), the required sets contain all of \(E\) eventually.

\hypertarget{1035-an-elementary-paduxe9-remainder-identity}{%
\section{4. An elementary Padé remainder
identity}\label{1035-an-elementary-paduxe9-remainder-identity}}

For every integer \(n\ge1\), there are real polynomials \(q_n,p_n\) such
that \(q_n\) is monic of degree \(n\), \(\deg p_n\le n-1\), and

\[
e^x q_n(x)-p_n(x)
=\frac{x^{2n}}{(n-1)!}
\int_0^1 e^{tx}t^n(1-t)^{n-1}\,dt.
\tag{1}
\]

Here is a direct proof. Let \(f_n(t)=t^n(1-t)^{n-1}\). Integrating by
parts repeatedly, using \(f_n^{(2n)}=0\), gives

\[
\frac{x^{2n}}{(n-1)!}\int_0^1 e^{tx} f_n(t)\,dt
=\frac1{(n-1)!}\sum_{j=0}^{2n-1}(-1)^j x^{2n-1-j}
\bigl(e^x f_n^{(j)}(1)-f_n^{(j)}(0)\bigr).
\]

The endpoint vanishing orders are \(n-1\) at \(1\) and \(n\) at \(0\).
Therefore the coefficient of \(e^x\) is a polynomial of degree \(n\),
while the other polynomial has degree at most \(n-1\). Since
\(f_n^{(n-1)}(1)=(-1)^{n-1}(n-1)!\), the former polynomial is monic.
This proves (1), initially for \(x\ne0\) and then at \(0\) by
continuity.

For reference, explicit coefficients are

\[
q_n(x)=\sum_{k=0}^{n}(-1)^{n+k}
\frac{(2n-1-k)!\,n!}{(n-1)!\,(n-k)!\,k!}x^k,
\]

\[
p_n(x)=(-1)^n\sum_{k=0}^{n-1}
\frac{(2n-1-k)!}{(n-1-k)!\,k!}x^k.
\]

These are the monically rescaled Padé approximants of type \([n-1/n]\)
for the exponential. No Padé approximation theorem is needed: identity
(1) is sufficient.

For \(|x|\le1\), the beta integral yields

\[
|e^xq_n(x)-p_n(x)|
\le\frac{e}{(n-1)!}\int_0^1t^n(1-t)^{n-1}\,dt
= e\frac{n!}{(2n)!}.
\tag{2}
\]

\hypertarget{1035-the-determinant-regularity-limit-is-zero}{%
\section{5. The determinant regularity limit is
zero}\label{1035-the-determinant-regularity-limit-is-zero}}

Let \(P_n\) denote the monic matrix orthogonal polynomial for \(\mu\),
and let

\[
H_n=\int P_n(x)^*\,d\mu(x)\,P_n(x).
\]

For every monic matrix polynomial \(Q\) of degree \(n\), orthogonality
gives

\[
\int Q^*\,d\mu\,Q
=H_n+\int(Q-P_n)^*\,d\mu\,(Q-P_n)\succeq H_n.
\tag{3}
\]

Consider the monic trial matrix

\[
T_n(x)=\begin{pmatrix}x^n&-p_n(x)\\0&q_n(x)\end{pmatrix}
\]

for \(\nu\), and \(\widetilde T_n=C^{-1}T_n C\) for \(\mu\). Both have
leading coefficient \(I_2\). If \(R_n=e^xq_n-p_n\), then

\[
K_n:=\int T_n^*\,d\nu\,T_n
=\int_{-1}^1
\begin{pmatrix}x^{2n}&x^nR_n(x)\\x^nR_n(x)&R_n(x)^2\end{pmatrix}\,dx.
\]

The determinant of this Gram matrix is bounded above by the product of
its diagonal entries. By (2),

\[
\det K_n\le
\left(\int_{-1}^1x^{2n}\,dx\right)
\left(\int_{-1}^1R_n(x)^2\,dx\right)
\le4e^2\left(\frac{n!}{(2n)!}\right)^2.
\tag{4}
\]

The normalized trial norm is \(CK_n C\). Since determinants are monotone
for the positive-definite order, (3)--(4) imply

\[
\det H_n\le\det(CK_nC)
\le\frac{4e^2}{\det M}\left(\frac{n!}{(2n)!}\right)^2.
\tag{5}
\]

For a normalized matrix measure, comparison of the leading coefficients
in the orthonormal three-term recurrence gives

\[
\det H_n=|\det(A_1\cdots A_n)|^2.
\tag{6}
\]

For completeness, the recurrence gives

\[
|\det A_j|=(\det H_j/\det H_{j-1})^{1/2}.
\]

Multiplying and using \(H_0=I_2\) gives (6). This is also \cite[Eq.~(2.34)]{DPS08}.

Combining (5) and (6),

\[
0<|\det(A_1\cdots A_n)|
\le \frac{2e}{\sqrt{\det M}}\frac{n!}{(2n)!}.
\]

Finally,

\[
\frac{n!}{(2n)!}
=\frac1{(n+1)(n+2)\cdots(2n)}\le(n+1)^{-n}.
\]

Taking \(n\)-th roots proves that the determinant limit is zero. This
contradicts regularity, whose required value is \(1/4\), and completes
the counterexample.

\hypertarget{1035-why-the-scalar-intuition-fails}{%
\section{Why the scalar intuition
fails}\label{1035-why-the-scalar-intuition-fails}}

Every fixed vector direction receives substantial mass on every short
interval. Nevertheless, the null direction of the rank-one density
rotates analytically with \(x\). Vector polynomials can follow this
changing null direction to superexponential accuracy, via rational
approximation of \(e^x\). The local integral lower bound does not
control such changing polynomial directions. This mechanism is absent
from a scalar measure or a fixed diagonal direct sum.

\begin{thebibliography}{1}
\raggedright
\bibitem{DPS08}
D. Damanik, A. Pushnitski and B. Simon,
\emph{The Analytic Theory of Matrix Orthogonal Polynomials},
Surveys in Approximation Theory \textbf{4} (2008), 1--85.
\url{https://www.math.auckland.ac.nz/sat/papers/11/11.pdf}.
\url{https://arxiv.org/abs/0711.2703v2}.
The positive-semidefinite convention is in Section~1.2 (p.~5),
polynomial nontriviality in Lemma~2.1 (p.~16),
the norm identity in (2.34) (p.~21), determinant regularity in
Theorem~5.4 (p.~70), and Conjecture~5.9 in (5.20) (p.~72).
\end{thebibliography}
\end{document}
